\section{Experiments}
We organize the experiments around the link between predicted difficulty and useful computation. The original two-dataset results establish the accuracy--graph tradeoff to be explained. We then hold histories and optional-pair budgets fixed to test the allocation signal, examine alternative explanations for its behavior, and finally return to autonomous rollouts to measure accumulated error and cost.

\subsection{The original benchmark tradeoff}
The original Sand and WaterDrop experiments cover granular and fluid dynamics, respectively, with 30 saved trajectories and one trained model per method in each comparison. Their contrasting behavior motivates a controlled study of allocation. On Sand, adaptation has lower full-horizon mean error than sparse MSE but uses more directed edges; the conditional error interval crosses zero. On WaterDrop, it uses fewer edges but has higher full-horizon error (Table~\ref{tab:historical-main}). Its MSE at forecast 200 improves by 0.327\%, with a conditional interval crossing zero, whereas mean error over the full rollout is 3.86\% worse. Against the fixed legacy risk-trained checkpoint, adaptive evaluation has lower MSE at forecast 200 on both datasets (Appendix~\ref{sec:historical-audit}).

\begin{table}[t]
\centering\small
\caption{Original saved-model experiments: adaptive minus sparse MSE. MSE is averaged over the full saved horizon; edges are mean directed counts. Negative differences favor adaptation. These are historical comparisons, not corrected retraining.}
\label{tab:historical-main}
\begin{tabular}{lrrr}
\toprule
Dataset & Steps & $\Delta$MSE & $\Delta$edges\\
\midrule
Sand & 314 & $-.0015016$ & $+4276.46$\\
WaterDrop & 995 & $+.0013357$ & $-1919.84$\\
\bottomrule
\end{tabular}
\end{table}

These comparisons leave two explanations entangled: the effect of extra messages on a prediction, and the effect of changed predictions on future geometry. The saved arrays also omit checkpoint hashes, training seeds and trajectory IDs, limiting verification of their training, model selection and row pairing. We treat their intervals as conditional on the saved models and assumed row alignment. The next experiments address the allocation question using matched model inputs and explicit budgets.

\subsection{Controlled models, histories, and policies}
We train six full-architecture WaterDrop models: width 128, ten message-passing blocks, and 100,000 updates on all 1,000 official training trajectories. Corrected Gaussian NLL and faithful regression each use three seeds, with paired initialization and frame/noise schedules. Faithful regression controls whether variance fitting changes the mean-training gradients; it is an established objective, not a new loss. Evaluation uses the fixed final checkpoints. The paired validation results show why this control matters: NLL has higher clean MSE in every seed despite smaller binned risk gaps (Appendix~\ref{sec:full-protocol}). A smaller binned risk gap need not accompany a more accurate mean predictor.

Base and dense graphs test the amount of expansion. Random25, speed25 and risk25 each retain exactly one quarter of the available annulus pairs, rounded down, so their common-state differences test placement. We retain the original radius ratio of 1.267. Observed-history comparisons use 297 histories per model from test trajectories 3--29; separate graph-convention controls also use 128 validation histories. Previous-observed risk is recomputed on the preceding observed base history. Autonomous cached risk instead reuses its own preceding selected-graph scores. This distinction separates the information in a score on a fixed input from its behavior under feedback.

The initial controlled evaluation uses uncapped symmetric pairs without self-loops, whereas training uses a capped base graph with self-loops and no annulus expansion. We test the self-loop and cap differences below; expanded-graph training support remains a separate limitation. The full-model protocol fixes checkpoints and policies before its evaluation, but the first three test trajectories and historical aggregates were previously inspected. Subsequent action-benefit, graph-convention and exposure analyses are exploratory. Exact protocols, data identities and implementation corrections are in the appendix.

A complementary five-seed compact study examines corrected NLL, beta-NLL and faithful regression, benefit prediction and inexpensive physical selectors. It again ranks residuals more strongly than dense benefit, but policy ordering varies by objective and split. Its architecture, training budget, noise and expanded-graph exposure differ from the full models, so the comparison does not isolate one cause of that variation. Appendix~\ref{sec:compact-study} retains the complete methods and Tables~\ref{tab:pilot} and~\ref{tab:rollout}; Appendices~\ref{sec:physical-rank} and~\ref{sec:physical-allocation} report the physical diagnostics.

\subsection{Does residual ranking identify useful messages?}
The controller's premise first requires a score that identifies difficult predictions. On the fixed observed histories, previous-base risk has mean frame correlations of $.357\pm.060$ (faithful) and $.411\pm.023$ (NLL) with base residuals. Its correlations with signed dense-intervention benefit are only $-.022\pm.045$ and $-.048\pm.012$. Random allocation also has lower one-step MSE than previous risk in every seed of both objectives. Scores therefore rank residual error without yielding the same ordering of useful graph changes.

Dense expansion is only a proxy for the controller's actual action: the policy selects 25\% of optional pairs, not all of them. To test whether that proxy explains the weak relationship, we use the saved predictions to measure the signed benefit of each selected sparse graph. Previous-risk correlation with its own action's normalized vector benefit is $-0.021\pm0.049$ for faithful and $-0.052\pm0.013$ for NLL (Figure~\ref{fig:actual-action-risk}). Dense and actual-risk benefit signs disagree for $30.2\pm0.6$\% and $31.3\pm0.5$\% of particles. The action definition matters, but measuring the chosen action does not recover strong risk ranking in these models.

\begin{figure}[t]
\centering\setlength{\unitlength}{1pt}
\begin{picture}(225,167)
\put(12,155){\makebox(0,0)[l]{\scriptsize Previous-risk Spearman correlation}}
\put(30,31){\line(1,0){181}}
\put(30,31){\line(0,1){114}}
\put(27,31.00){\makebox(0,0)[r]{\tiny -0.1}}
\put(29,31.00){\line(1,0){3}}
\put(27,49.00){\makebox(0,0)[r]{\tiny 0.0}}
\put(29,49.00){\line(1,0){3}}
\put(27,67.00){\makebox(0,0)[r]{\tiny 0.1}}
\put(29,67.00){\line(1,0){3}}
\put(27,85.00){\makebox(0,0)[r]{\tiny 0.2}}
\put(29,85.00){\line(1,0){3}}
\put(27,103.00){\makebox(0,0)[r]{\tiny 0.3}}
\put(29,103.00){\line(1,0){3}}
\put(27,121.00){\makebox(0,0)[r]{\tiny 0.4}}
\put(29,121.00){\line(1,0){3}}
\put(27,139.00){\makebox(0,0)[r]{\tiny 0.5}}
\put(29,139.00){\line(1,0){3}}
\multiput(32,49)(4,0){45}{\line(1,0){1}}
\put(50.00,101.43){\circle*{3.5}}
\put(52.00,122.47){\circle*{3.5}}
\put(54.00,115.92){\circle*{3.5}}
\put(48.00,113.27){\line(1,0){8}}
\put(111.00,45.20){\circle*{3.5}}
\put(113.00,53.06){\circle*{3.5}}
\put(115.00,36.93){\circle*{3.5}}
\put(109.00,45.06){\line(1,0){8}}
\put(172.00,47.70){\circle*{3.5}}
\put(174.00,52.41){\circle*{3.5}}
\put(176.00,35.43){\circle*{3.5}}
\put(170.00,45.18){\line(1,0){8}}
\put(64.00,120.05){\makebox(0,0){\tiny$\triangle$}}
\put(66.00,127.68){\makebox(0,0){\tiny$\triangle$}}
\put(68.00,121.38){\makebox(0,0){\tiny$\triangle$}}
\put(62.00,123.04){\line(1,0){8}}
\put(125.00,42.54){\makebox(0,0){\tiny$\triangle$}}
\put(127.00,38.39){\makebox(0,0){\tiny$\triangle$}}
\put(129.00,40.05){\makebox(0,0){\tiny$\triangle$}}
\put(123.00,40.33){\line(1,0){8}}
\put(186.00,42.11){\makebox(0,0){\tiny$\triangle$}}
\put(188.00,39.52){\makebox(0,0){\tiny$\triangle$}}
\put(190.00,37.41){\makebox(0,0){\tiny$\triangle$}}
\put(184.00,39.68){\line(1,0){8}}
\put(59,21){\makebox(0,0){\tiny Base residual}}
\put(120,21){\makebox(0,0){\tiny Dense benefit}}
\put(181,21){\makebox(0,0){\tiny Actual risk25}}
\put(68,7){\circle*{3.5}}
\put(74,7){\makebox(0,0)[l]{\tiny Faithful}}
\put(136,7){\makebox(0,0){\tiny$\triangle$}}
\put(143,7){\makebox(0,0)[l]{\tiny NLL}}
\end{picture}
\caption{Residual ranking and actual action value on the saved observed histories. Each symbol is one seed mean; horizontal marks are three-seed means. Actual risk25 benefit uses the whole selected graph. This exploratory comparison retains the original no-self-loop evaluation convention.}
\label{fig:actual-action-risk}
\end{figure}

Weak mean rank association can coexist with useful behavior in part of the distribution. The highest risk quartile has positive mean risk-action benefit for each objective, although one NLL seed is negative. Random allocation nevertheless provides larger benefit in that quartile in all six models. A hindsight choice among base, random25, speed25 and risk25 selects base as strictly best on $35.9\pm12.0$\% / $35.8\pm6.3$\% of frames (faithful/NLL). This target-using reference suggests that deciding whether to expand can matter alongside deciding where to expand; it is not a deployable selector. We average frames within each trajectory and weight trajectories equally within each seed, then report means and sample SDs across the three seeds.

\subsection{Does the graph convention explain the deficit?}
Removing trained self-messages could affect both residual prediction and the response to extra edges. We therefore compare the same six checkpoints on 425 observed histories each (128 validation, 297 test), crossing self-loops and the 128-neighbor cap for base/dense graphs and testing both loop settings for every selector. Restoring self-loops lowers base test MSE by $17.85\pm0.75$\% / $17.21\pm3.01$\% (faithful/NLL; paired within-seed changes).

Absolute accuracy changes substantially, yet previous-observed risk remains worse than random, and dense worse than base, in every seed on both splits. Fresh current-base risk also fails to beat random. With self-loops, the residual--benefit separation also persists: previous-base risk has three-seed mean test Spearman correlations $0.325/0.386$ with base residuals versus $-0.030/-0.075$ with its own sparse-action benefit (faithful/NLL). Faithful test benefit correlations change sign across seeds; Table~\ref{tab:graph-bridge-risk-benefit} gives both splits and risk rules. The cap is inactive on these observed histories, and feature/prediction checks establish agreement with the native base graph. The result survives restoration of trained self-messages and fresh scoring; it does not establish robustness to expanded-graph training. The autonomous control below tests the remaining feedback question. Appendix~\ref{sec:graph-bridge} reports every policy and the paired loop interactions.

\subsection{How does allocation change the prediction?}
The preceding controls leave a fixed-state risk--random gap to explain. Returning to the original no-self-loop predictions, we decompose it by how the graph changes the model output. Risk-minus-random normalized coordinate MSE, multiplied by $10^3$, is $0.486\pm0.111$ / $0.860\pm0.102$ (faithful/NLL). Let $r_i$ be the normalized truth-minus-base residual and $\delta_i$ the normalized change from the base prediction. With particle means $A=\langle2r_i^\top\delta_i/d\rangle_i$ and $C=\langle\|\delta_i\|^2/d\rangle_i$ ($d=2$), risk-minus-random error obeys $\Delta E=\Delta C-\Delta A$. The positive gap means that the larger prediction change outweighs the alignment gain; $\Delta A>0$ in every seed is the additional empirical observation. Here $C$ denotes squared prediction change, not computational cost.

We also partition particles by whether neither policy, only risk, only random, or both policies deliver optional messages directly to them. The both-covered group makes the largest positive unconditional contribution in every seed, so the deficit is not confined to particles reached by only one selector. These policy-defined groups describe where the error appears; they do not identify a causal concentration effect. Ten message-passing blocks also allow indirect effects on particles in the neither group. Appendix~\ref{sec:optional-exposure} retains every group's size and contribution and both speed controls.

\subsection{Does the allocation help over a full rollout?}
A useful one-step action must also withstand feedback: each prediction changes the positions and graphs seen at later steps. We therefore evaluate all five policies on the same 27 test trajectories for 995 forecasts, observing only the initial six frames. Table~\ref{tab:full-rollouts} reports full-horizon position MSE and failures. A required failed trajectory leaves the all-sample full-horizon group mean undefined; accepted-prefix and boundary diagnostics are reported separately.

\begin{table*}[t]
\centering\small
\caption{Full-architecture position MSE over 995 forecasts, mean $\pm$ sample SD of three equal-trajectory seed means. Failures are counts out of 81 outcomes per objective/policy, shown faithful / NLL. A failed trajectory makes the full-horizon group mean undefined.}
\label{tab:full-rollouts}
\begin{tabular}{lrrr}
\toprule
Policy & Faithful & NLL & Failures\\
\midrule
Base & $0.0266\pm0.0202$ & --- & 0 / 1\\
Dense & $0.0382\pm0.0188$ & --- & 0 / 2\\
Random25 & $0.0274\pm0.0201$ & --- & 0 / 4\\
Speed25 & $0.0252\pm0.0148$ & $0.0321\pm0.0141$ & 0 / 0\\
Cached risk25 & $0.0282\pm0.0179$ & --- & 0 / 1\\
\bottomrule
\end{tabular}
\end{table*}

All 405 faithful outcomes complete. Cached risk has higher faithful full-rollout MSE than speed in every seed, with paired difference $0.00301\pm0.00315$. Differences from base and random are $0.00155\pm0.00263$ and $0.00083\pm0.00225$ and change sign across seeds; cached risk improves on dense in all three seeds. NLL has eight coordinate-guard failures, all at seed 2, making its cached-risk full-horizon mean undefined. These results do not support a consistent autonomous advantage for cached risk.

A further native-convention rollout control retains all 810 outcomes (Appendix~\ref{sec:native-rollout-followup}). Faithful cached risk loses to base in all three seeds, with paired MSE difference $.004287\pm.003358$; its random comparison changes sign ($.000447\pm.000721$). Seven NLL guard failures leave all three-seed NLL full-horizon means undefined. Matching the trained base convention therefore does not resolve the allocation deficit.

Numerical completion is also weaker than physical plausibility. Faithful cached risk places $13.11\pm6.10\%$ of particles outside the metadata bounds by more than $10^{-6}$, versus $2.68\%$ for the same ground-truth frames. Its mean trajectory-maximum excursion is $.302$, compared with $.239$ for base and $.00782$ for truth. The nonzero truth reference matters: a box crossing alone is not an exact physical-failure label. Full endpoint, failure and boundary results are in Appendix~\ref{sec:full-results}.

\subsection{What computation does the policy require?}
Finally, the allocation benefit must be weighed against graph construction, selection and scoring. Six rotated, synchronized timing repeats on each observed history give natural-base costs of 10.74/10.50 ms and random costs of 11.80/11.55 ms (faithful/NLL). Previous-observed risk costs 22.38/21.88 ms because this diagnostic regenerates the preceding base score. Autonomous cached risk instead reuses its previous forward pass, so these common-state measurements do not measure its steady-state deployment cost. Every policy executes the risk head. Together with the geometry decomposition, the timings explain why neither lower rollout edge counts nor cached-score reuse alone establishes an accuracy--runtime advantage.

